% Nouveau EGA V (cote 25): archive cote 25, site page 53 (PDF page 54).
% typescript. Draft compared: none.
% AI-transcribed by Claude Opus 5.5 (claude-opus-5-5, adaptive thinking, effort high), Message Batches API, batch ega25b (msgbatch_01A7n5YuXubUPNK1zsepi8RG), 2026-11-13; not yet checked by eye.
\scanpage{25/53}
\var{\margin{\hand{\unsure{46|48}}}}{\placedmargin{\hand{\unsure{46|48}}}}{marginal note placed}{ega25_25_053-v0001}

$\overset{\hand{\vee}}{P}$ (correspondant à la question d'étudier les sections hypérplanes de X
passant par un point donné x, ou tangent à X en un point \var{\deleted{simple}}{}{deletion}{ega25_25_053-v0002}\var{\inserted{lisse}}{lisse}{insertion}{ega25_25_053-v0003} donné),
form\var{\over{e}{é}}{é}{writing over}{ega25_25_053-v0004} donc des hyperplans contenant une sous-variété linéaire $L^{o}$ de P
(resp. un point,  ou l'image d'un espace tangent de X en un point lisse,
dans les deux cas envisagés), et on se pose la question de savoir si
pour le point générique \hand{$\xi$} de L  (donc pour tous les points d'un ouvert
non vide de L),  $Y_{\hand{\xi}}$  est \var{\xout{géom intègre ou du moins}}{}{typewriter strike-out}{ega25_25_053-v0005} géom irréd., de
dim  dim X -1.   C'est là une variante  de Bertini, qui devrait figurer
dans  le n°\hand{3}  et se traite par exactement la même méthode\var{\inserted[typed]{ou, si on veut, s'y ramène.}}{ ou, si on veut, s'y ramène.}{insertion}{ega25_25_053-v0006}  La question
dimensionnelle s'exprime simplement par  $f(X)\hand{\not\subset} L^{o}$ , \var{\xout{Il y a donc un}}{}{typewriter strike-out}{ega25_25_053-v0007}
i.e.  $X' =$ \var{\xout{X=}}{}{typewriter strike-out}{ega25_25_053-v0008} $f^{-1}(P - L^{o})$   est un ouvert dense de X, \var{\xout{et la question}}{}{typewriter strike-out}{ega25_25_053-v0009}
Soit \unsure{Q|N} l'espace projectif des hypérplans passant par $L^{o}$  (NB si  $L^{o}$ est
défini par un sous-espace vectoriel $F^{o}$ de E,  on a   $Q = P($\var{\deleted{E/}}{}{deletion}{ega25_25_053-v0010}$F^{o})$ ) et
considérons le morphisme  canonique  (déduit de \var{\xout{E\illegible{1}E/$F^{o}$}}{}{typewriter strike-out}{ega25_25_053-v0011} $F^{o}\hand{\longrightarrow} E$ , cf
Chap II)
\[ u :\ P - L^{o} \hand{\longrightarrow} Q, \]
et considérons
\[ g = u f' :\quad f^{-1}(P - L^{o}) = X' \hand{\longrightarrow} Q , \]
alors \var{\xout{la famille des X' \illegible{1} n'est autr}}{}{typewriter strike-out}{ega25_25_053-v0012}  $L \hand{\simeq} \overset{\hand{\vee}}{Q}$ ,  et la famille des
$X'_{\hand{\xi}}$  $(\hand{\xi}\hand{\in} L)$  n'est autre que la famille des ``sections hyperplanes'' rela-%
tif à ce morphisme  g . D'autre part, on voit aussitôt que pour tout
$\hand{\xi}\hand{\in} L$  ``général''  $X'_{\hand{\xi}}$ est dense dans $X_{\hand{\xi}}$ , donc  $X'_{\hand{\xi}}$  est géom irréd sss
X  l'est.  Ceci posé,  le théorè\unsure{m|p}e de Bertini-Zariski nous montre que
nous avons la conclusion voulue \var{\xout{dix}}{}{typewriter strike-out}{ega25_25_053-v0013} d'irréductibilité pourvu que
dim $g(X') \hand{\geqslant} 2$. (A vrai dire, tant qu'à faire on pourrait signaler une
réciproque dans \hand{3}.1. , savoir : \var{\xout{Y est gé}}{}{typewriter strike-out}{ega25_25_053-v0014} si X est géom irréd, Y  est
géom irréd sss  ou bien  dim $f(X)$; $\hand{\neq}$ \unsure{2|1}, ou bien  dim $f(X)=1$ , et $f(X)$
est contenu dans une droite\hand{D} (définie sur $\overset{\hand{\rule{0.6em}{0.35pt}}}{k}$ )\var{\inserted{et la fibre générique de $X\to D$ est géom. \unsure{irréd}}}{ et la fibre générique de $X\to D$ est géom. \unsure{irréd}}{insertion}{ega25_25_053-v0015}\hand{)}; cela permet également

\var{\margin{\hand{53}}}{\placedmargin{\hand{53}}}{marginal note placed}{ega25_25_053-v0016}

% NOTE: Beside dim dim X-1 the left margin has several parallel pencil strokes; no textual slash is supplied.
